% ============================================================================
\section{Additional Mitigation Controls}
\label{app:mitigations}

The main text focuses on the interventions that directly test the readout-blind-spot mechanism.
We also ran mitigation-style controls that help position the result relative to existing looped-transformer designs.

\paragraph{Exit-weighted objective.}
An Ouro-style exit-probability-weighted objective keeps RMSNorm readouts but changes how per-loop CE terms are weighted.
It reduces neither the qualitative drift pattern nor the radial blind spot: final-loop norms remain in the thousands at both scales ($6{,}470 \pm 2{,}233$ at 44M and $4{,}710 \pm 725$ at 129M).
This supports the distinction between changing CE weights and changing what CE can see.

\paragraph{Spectral damping.}
Spectral damping multiplies the hidden state by a learned factor below 1 between loops.
It limits amplification but does not make the readout loss norm-sensitive.
In our runs it lowers norms relative to terminal + RMSNorm ($2{,}432 \pm 238$ at 44M and $2{,}973 \pm 16$ at 129M on the stable seeds), but remains far above raw-readout norms.
One of three 129M seeds became unstable, with hidden-state norms exceeding $10^5$, so we treat it as less robust than exposing scale through the loss.

\paragraph{Inter-loop normalization.}
Ouro uses RMSNorm between loop applications.
This is a different solution: it caps the scale entering each recurrent update rather than making the output loss scale-sensitive.
It is effective, but it changes the loop dynamics and adds normalization inside the recurrent path.

\paragraph{Latent-space losses.}
In denoising-style looped experiments, per-loop latent losses can prevent drift because they directly train $H_k$ in representation space.
This is consistent with our mechanism.
Per-loop losses help scale only when the loss itself is scale-sensitive; CE through a normalized readout is not.


